\chapter{Kombinatorika dan Pencacahan}\label{ch:g12:comb}

Kombinatorika adalah seni mencacah tanpa mendaftar. Dua asasnya yang paling
dasar --- jumlahkan ukuran pilihan yang \omterm{def:g10:proba:operations}{saling lepas}, kalikan banyaknya
pilihan yang saling bebas --- sudah cukup untuk mencacah \omterm{def:g12:comb:tuples}{susunan}, \omterm{def:g12:comb:permutation}{permutasi}
dan himpunan bagian suatu himpunan berhingga, dan berpuncak pada teorema binomial.

\section{Dua asas pencacahan}

Kita menulis $\abs{E}$ untuk banyaknya anggota (\emph{kardinalitas}) suatu
himpunan berhingga $E$.

\begin{proposition}[Asas penjumlahan]\label{prop:g12:comb:addition}
Jika himpunan berhingga $E$ dipartisi menjadi himpunan bagian $A_1, \dots, A_k$
(\omterm{def:g10:proba:operations}{saling lepas} sepasang demi sepasang, dengan gabungan $E$), maka
\[
\abs{E} = \abs{A_1} + \abs{A_2} + \dots + \abs{A_k}.
\]
\end{proposition}

\begin{proposition}[Asas perkalian]\label{prop:g12:comb:multiplication}
Jika sebuah objek dibangun lewat $k$ pilihan berturut-turut, dengan $n_1$ pilihan
pada langkah pertama dan, \emph{apa pun pilihan sebelumnya}, $n_i$ pilihan pada
langkah ke-$i$, maka banyaknya objek yang terbangun adalah
$n_1 \times n_2 \times \dots \times n_k$.
\end{proposition}

\begin{proof}
Kedua pernyataan itu dibuktikan dengan induksi pada $k$; kasus $k = 2$ pada
pernyataan kedua sama dengan mencacah larik persegi panjang baris demi baris.
\end{proof}

\begin{example}
Sebuah rumah makan menawarkan 4 hidangan pembuka, 6 hidangan utama, 3 hidangan
penutup: $4 \times 6 \times 3 = 72$ santapan tiga hidangan yang berbeda.
\end{example}

\section{Tupel, permutasi, faktorial}

\begin{definition}[Tupel-$k$]\label{def:g12:comb:tuples}
Suatu \emph{tupel-$k$} dari himpunan $E$ adalah daftar terurut $(x_1, \dots, x_k)$
beranggotakan $E$, dengan pengulangan diperbolehkan. Tupel-$k$ yang anggotanya
\emph{berbeda} disebut \emph{susunan}\index{susunan} $k$ anggota $E$.
\end{definition}

\begin{proposition}\label{prop:g12:comb:tuples}
Misalkan $\abs E = n$. Banyaknya tupel-$k$ dari $E$ adalah $n^k$. Banyaknya
\omterm{def:g12:comb:tuples}{susunan} $k$ anggota $E$ ($0 \leq k \leq n$) adalah
\[
n(n-1)(n-2)\cdots(n-k+1) = \frac{n!}{(n-k)!},
\]
dengan $n! = 1 \times 2 \times \dots \times n$ (dan $0! = 1$) yang disebut
\emph{faktorial}\index{faktorial} dari $n$.
\end{proposition}

\begin{proof}
Asas perkalian: untuk tupel-$k$ ada $n$ pilihan pada masing-masing
$k$ langkahnya; untuk \omterm{def:g12:comb:tuples}{susunan}, ada $n$ pilihan bagi $x_1$, lalu $n - 1$ bagi
$x_2$ (satu anggotanya sudah terpakai), \dots, $n - k + 1$ bagi $x_k$.
\end{proof}

\begin{definition}[Permutasi]\label{def:g12:comb:permutation}
Suatu \emph{permutasi}\index{permutasi} dari $E$ adalah \omterm{def:g12:comb:tuples}{susunan} seluruh $n$
anggota $E$, yaitu pengurutan $E$. Menurut
\cref{prop:g12:comb:tuples} (kasus $k = n$), banyaknya permutasi suatu
himpunan beranggota $n$ adalah $n!$.
\end{definition}

\begin{example}
Lima pelari dapat mencapai garis akhir dalam $5! = 120$ urutan yang berbeda.
Banyaknya podium yang mungkin (tiga tempat teratas) adalah $5 \times 4 \times 3 = 60$.
\end{example}

\section{Kombinasi dan koefisien binomial}

\begin{definition}[Kombinasi]\label{def:g12:comb:combination}
Suatu \emph{kombinasi}\index{kombinasi} dari $k$ anggota $E$ adalah himpunan
bagian $E$ yang beranggota $k$ (tanpa urutan, tanpa pengulangan). Banyaknya ditulis
$\dbinom{n}{k}$\index{koefisien binomial}, dibaca ``$n$ pilih $k$''.
\end{definition}

\begin{theorem}\label{thm:g12:comb:binomformula}
Untuk $0 \leq k \leq n$:
\[
\binom{n}{k} = \frac{n!}{k!\,(n-k)!} .
\]
\end{theorem}

\begin{proof}
Cacahlah \omterm{def:g12:comb:tuples}{susunan} $k$ anggota $E$ dengan dua cara. Secara langsung:
$\frac{n!}{(n-k)!}$. Atau, pilihlah dahulu himpunan bagian yang mendasarinya
($\binom nk$ cara), lalu urutkan ($k!$ cara); asas perkaliannya
memberi $\binom{n}{k}\,k!$. Dengan menyamakan keduanya,
$\binom nk = \frac{n!}{k!(n-k)!}$.
\end{proof}

\begin{proposition}[Identitas dasar]\label{prop:g12:comb:identities}
Untuk $0 \leq k \leq n$:
\[
\binom{n}{0} = \binom{n}{n} = 1, \qquad
\binom{n}{1} = n, \qquad
\binom{n}{k} = \binom{n}{n-k},
\]
dan \emph{kaidah Pascal}\index{segitiga Pascal}: untuk $1 \leq k \leq n-1$,
\[
\binom{n}{k} = \binom{n-1}{k-1} + \binom{n-1}{k}.
\]
\end{proposition}

\begin{proof}
Kesetangkupan $\binom nk = \binom{n}{n-k}$ berlaku sebab pengambilan \omterm{def:g10:proba:operations}{komplemennya}
memasangkan himpunan bagian beranggota $k$ dengan yang beranggota $(n-k)$, satu
lawan satu. Untuk \omterm{prop:g12:comb:identities}{kaidah Pascal}, tetapkan satu anggota $a \in E$ lalu pilahlah
himpunan bagian beranggota $k$ itu menjadi yang memuat $a$ --- yang diperoleh
dengan menambahkan $a$ pada himpunan bagian beranggota $(k-1)$ dari
$E \setminus \{a\}$, dan banyaknya $\binom{n-1}{k-1}$ --- serta
yang tidak memuat $a$, yaitu himpunan bagian beranggota $k$ dari
$E \setminus \{a\}$, yang banyaknya $\binom{n-1}{k}$. Simpulkan dengan asas
penjumlahannya.
\end{proof}

\omterm{prop:g12:comb:identities}{Kaidah Pascal} membangkitkan koefisiennya baris demi baris, dan itulah
yang disebut \emph{\omterm{prop:g11:binom:pascal}{segitiga Pascal}}: setiap isian adalah jumlah dua isian di atasnya.

\begin{omfigure}
\begin{tikzpicture}[x=0.95cm, y=0.8cm, font=\small]
  \node at (0,0) {$1$};
  \node at (-0.5,-1) {$1$};   \node at (0.5,-1) {$1$};
  \node at (-1,-2) {$1$};     \node at (0,-2) {$2$};
  \node at (1,-2) {$1$};
  \node at (-1.5,-3) {$1$};   \node at (-0.5,-3) {$3$};
  \node at (0.5,-3) {$3$};    \node at (1.5,-3) {$1$};
  \node at (-2,-4) {$1$};
  \node[circle, draw=omProp, inner sep=1.5pt] (p1) at (-1,-4) {$4$};
  \node[circle, draw=omProp, inner sep=1.5pt] (p2) at (0,-4) {$6$};
  \node at (1,-4) {$4$};      \node at (2,-4) {$1$};
  \node at (-2.5,-5) {$1$};   \node at (-1.5,-5) {$5$};
  \node[circle, draw=omThm, inner sep=1.5pt] (s) at (-0.5,-5) {$10$};
  \node at (0.5,-5) {$10$};   \node at (1.5,-5) {$5$};
  \node at (2.5,-5) {$1$};
  \draw[-Stealth, omProp] (p1) -- (s);
  \draw[-Stealth, omProp] (p2) -- (s);
\end{tikzpicture}

{\small \omterm{prop:g11:binom:pascal}{Segitiga Pascal}, baris $n = 0$ sampai $5$: \omterm{prop:g12:comb:identities}{kaidah Pascal}
$\binom{4}{1} + \binom{4}{2} = \binom{5}{2}$ sedang beraksi.}
\end{omfigure}

\begin{theorem}[Teorema binomial]\label{thm:g12:comb:binomial}
\index{teorema binomial}
Untuk semua $a, b \in \R$ (atau $\C$) dan $n \in \N$:
\[
(a+b)^n = \sum_{k=0}^{n} \binom{n}{k}\, a^{k}\, b^{\,n-k} .
\]
\end{theorem}

\begin{proof}
Jabarkan hasil kali $(a+b)(a+b)\cdots(a+b)$ ($n$ faktor): setiap suku pada
penjabarannya memilih $a$ atau $b$ pada tiap faktornya, dan menghasilkan
$a^k b^{n-k}$ dengan $k$ menyatakan banyaknya faktor yang menyumbang $a$.
Banyaknya cara memilih $k$ faktor itu di antara $n$ faktornya adalah
$\binom nk$, dan itulah koefisien $a^k b^{n-k}$.
\end{proof}

\begin{corollary}\label{cor:g12:comb:sums}
$\displaystyle\sum_{k=0}^{n} \binom{n}{k} = 2^n$
\quad dan \quad
$\displaystyle\sum_{k=0}^{n} (-1)^k\binom{n}{k} = 0$ \ ($n \geq 1$).
\end{corollary}

\begin{proof}
Ambil $a = b = 1$, lalu $a = -1$, $b = 1$ pada teorema binomialnya. Identitas
pertamanya juga bermakna langsung: himpunan beranggota $n$ mempunyai $2^n$
himpunan bagian (tiap anggotanya masuk atau tidak: asas perkalian), yang
dipilah menurut ukurannya.
\end{proof}

\begin{method}[Memilih model yang tepat]\label{met:g12:comb:model}
Sebelum mencacah, jawablah dua pertanyaan: \emph{apakah urutannya penting?} dan
\emph{apakah pengulangan diperbolehkan?}
\begin{center}
\begin{tabular}{l|cc}
 & urutan penting & urutan tak penting\\
\hline
pengulangan boleh & $n^k$ (tupel) & (universitas)\\
tanpa pengulangan & $\frac{n!}{(n-k)!}$ (\omterm{def:g12:comb:tuples}{susunan}) & $\binom nk$ (himpunan bagian)
\end{tabular}
\end{center}
Mengambil bola dari sebuah guci: \emph{dengan pengembalian, berurutan} $\to$ tupel;
\emph{tanpa pengembalian, berurutan} $\to$ \omterm{def:g12:comb:tuples}{susunan}; \emph{segenggam sekaligus}
$\to$ himpunan bagian.
\end{method}

\section{Latihan}

\begin{exercise}[$\star$]\label{exo:g12:comb:1}
Sebuah pelat nomor terdiri atas 2 huruf (A--Z), lalu 3 angka, lalu 2 huruf.
Berapa banyak pelat yang mungkin? Berapa yang tak punya karakter berulang?
\end{exercise}

\begin{exercise}[$\star$]\label{exo:g12:comb:2}
Hitunglah $\dbinom{8}{3}$, $\dbinom{10}{8}$, lalu sederhanakan
$\dfrac{\binom{n}{2}}{\binom{n+1}{2}}$.
\end{exercise}

\begin{exercise}[$\star$]\label{exo:g12:comb:3}
Di sebuah kelas berisi 30 murid, harus dipilih satu panitia beranggota 4 murid,
lalu seorang ketua dan seorang bendahara dari dalam panitia itu (satu orang tidak
boleh merangkap kedua jabatan). Berapa banyak hasil yang mungkin?
\end{exercise}

\begin{exercise}[$\star$]\label{exo:g12:comb:4}
Jabarkan $(x + 2)^5$ dan $(1 - x)^6$ dengan teorema binomial. Berapakah
koefisien $x^3$ pada $(2x + 3)^7$?
\end{exercise}

\begin{exercise}[$\star\star$]\label{exo:g12:comb:5}
Satu tangan poker baku terdiri atas 5 kartu dari satu pak berisi 52 kartu.
\begin{enumerate}
  \item Ada berapa banyak tangan yang mungkin?
  \item Berapa banyak tangan yang memuat tepat satu kartu raja? Sekurang-kurangnya
        satu kartu raja?
  \item Berapa banyak tangan yang berbentuk tiga-dua (tiga kartu berperingkat
        sama, dua kartu berperingkat lain yang sama)?
\end{enumerate}
\end{exercise}

\begin{exercise}[$\star\star$]\label{exo:g12:comb:6}
Ada berapa banyak anagram (penyusunan ulang hurufnya, bermakna atau tidak) yang
dimiliki kata MATH? Dan kata BANANA?
(Petunjuk untuk BANANA: tempatkan dahulu ketiga huruf A-nya.)
\end{exercise}

\begin{exercise}[$\star\star$]\label{exo:g12:comb:7}
Buktikan identitas $k\dbinom{n}{k} = n\dbinom{n-1}{k-1}$ ($1 \leq k \leq n$)
dengan dua cara: lewat rumus \omterm{prop:g12:comb:tuples}{faktorialnya}, dan dengan mencacah lewat dua cara
pasangan (panitia beranggota $k$ orang, ketuanya) yang dipilih dari $n$ orang.
\end{exercise}

\begin{exercise}[$\star\star$]\label{exo:g12:comb:8}
Sebuah lintasan pada bidang berjalan dari $(0,0)$ ke $(m, n)$ dengan langkah
satuan ke timur atau ke utara. Tunjukkan bahwa banyaknya lintasan itu $\dbinom{m+n}{m}$.
\end{exercise}

\begin{exercise}[$\star\star\star$]\label{exo:g12:comb:9}
Buktikan \emph{identitas Vandermonde}: untuk $0 \leq k \leq m + n$,
\[
\binom{m+n}{k} = \sum_{j=0}^{k} \binom{m}{j}\binom{n}{k-j},
\]
dengan mencacah himpunan bagian beranggota $k$ dari himpunan yang terbelah menjadi
satu kelompok beranggota $m$ dan satu kelompok beranggota $n$. Simpulkan bahwa
$\displaystyle\sum_{j=0}^{n}\binom{n}{j}^{\!2} = \binom{2n}{n}$.
\end{exercise}

\begin{exercise}[$\star\star\star$]\label{exo:g12:comb:10}
Dengan teorema binomial, tunjukkan bahwa untuk semua $n \geq 1$,
\[
\sum_{k=1}^{n} k \binom{n}{k} = n\,2^{n-1}.
\]
(Petunjuk: turunkan $(1+x)^n$, atau pakailah \cref{exo:g12:comb:7}.)
\end{exercise}

\section{Soal: Seni mencacah dua kali}

\begin{problem}[{Soal akhir pekan --- bintang dan batang, topi yang
tertukar semua, dan identitas yang terbukti dengan mencacah satu hal dua cara}]
\label{pb:g12:comb:1}
Muslihat terdalam dalam kombinatorika sesungguhnya sederhana sekali: cacahlah
kumpulan yang sama dua kali, dengan dua cara yang berbeda, lalu samakan
jawabannya. Soal ini melatih model pada
\cref{met:g12:comb:model}, menambahkan satu teknik yang tak diperlukan
uraian bab ini --- \emph{bintang dan batang} milik pencacahan es krim
--- lalu mencacah \emph{topi yang tertukar semua} yang masyhur itu secara eksak, dan
menemukan bilangan $\frac1\eu$ menunggu di dasar tumpukan
topinya, kemunculannya yang ketiga di dalam buku ini.

\medskip
\noindent\textbf{Bagian I --- Memilih modelnya.}
\begin{enumerate}
  \item Cacahlah pelat nomor yang tersusun atas $2$ huruf lalu
        $3$ angka; kemudian anagram kata BANANA.
  \item Dari satu pak berisi $32$ kartu, cacahlah tangan beranggota $5$ kartu;
        lalu tangan yang memuat tepat $2$ dari $4$ kartu raja.
  \item Sebuah robot berjalan dari $(0,0)$ ke $(4,3)$ hanya dengan langkah
        satuan ke kanan atau ke atas: ada berapa lintasan? (Sandikan lintasannya
        sebagai kata dalam huruf R dan U.)
  \item Jabarkan $(1 + x)^4$ dengan teorema binomial
        (\cref{thm:g12:comb:binomial}); lalu hitung nilainya di
        $x = 1$ dan $x = -1$: dua identitas mana tentang
        bilangan $\binom nk$ yang jatuh dari situ?
  \item Buktikan dengan pencacahan ganda bahwa
        $k\binom nk = n\binom{n-1}{k-1}$ (cacahlah
        panitia-beserta-ketuanya dengan dua cara), lalu simpulkan
        $\sum_{k=0}^{n} k\binom nk = n\,2^{n-1}$.
\end{enumerate}

\medskip
\noindent\textbf{Bagian II --- Bintang dan batang.}
\begin{enumerate}[resume]
  \item Sebuah kedai es krim menjual $4$ rasa; kamu memesan $10$
        sendok (rasanya boleh berulang, urutan di dalam cangkirnya
        tidak penting). Sandikan pesanannya sebagai deretan $10$ bintang
        (sendoknya) yang dipisahkan oleh $3$ batang (pergantian rasanya), lalu
        cacahlah pesanannya.
  \item Cacahlah tripel \omterm{def:g10:numbers:sets}{bilangan bulat} tak negatif dengan
        $x + y + z = 12$.
  \item Cacahlah tripel \omterm{def:g10:numbers:sets}{bilangan bulat} \emph{positif} dengan
        $x + y + z = 12$ (substitusikan $x = 1 + x'$, dan seterusnya).
  \item Ada berapa monomial berbeda yang muncul pada penjabaran
        $(a + b + c)^5$?
  \item Ujilah kewarasan caranya: cacahlah pesanan $3$
        sendok dari $2$ rasa dengan rumusnya, lalu daftarkan
        semuanya dan bandingkan.
  \item Sebutkan dengan tepat di mana ``sendoknya identik'' masuk ke dalam
        penyandiannya --- lalu cacahlah apa yang terjadi jika sebaliknya
        sendoknya dimakan berurutan (kedudukannya berbeda), dengan
        daftar periksa pada \cref{met:g12:comb:model}.
\end{enumerate}

\medskip
\noindent\textbf{Bagian III --- Topi yang tertukar semua.}
Suatu \emph{\omterm{def:g12:comb:permutation}{permutasi} kacau} adalah pembagian ulang $n$ topi kepada
$n$ pemiliknya yang membuat \emph{tak seorang pun} menerima topinya sendiri;
misalkan $D_n$ mencacahnya. (\Cref{pb:g11:prob:1} menunjukkan bahwa satu tamu
rata-rata memperoleh kembali topinya sendiri --- kini kita mencacah pesta yang
sepenuhnya sial itu secara eksak.)
\begin{enumerate}[resume]
  \item Hitunglah $D_1$, $D_2$, $D_3$ dengan mendaftar, lalu $D_4$
        dengan sabar (atau dengan cerdik).
  \item Berilah alasan bagi rekurensi
        $D_n = (n - 1)\left(D_{n-1} + D_{n-2}\right)$: tamu 1
        menerima suatu topi $k \neq 1$ ($n - 1$ pilihan); lalu pilahlah
        menurut apakah tamu $k$ menerima topi 1 atau tidak.
        Periksalah bahwa itu menghasilkan kembali $D_4$, lalu hitunglah $D_5$.
  \item Untuk $n = 3$, buktikan dengan pemuatan--pengeluaran (kurangkan
        pemberian yang menetapkan sekurang-kurangnya satu topi, lalu tambahkan
        kembali kelebihan cacahnya) bahwa
        $D_3 = 3!\left(1 - \frac{1}{1!} + \frac{1}{2!} -
        \frac{1}{3!}\right)$, lalu nyatakan rumus umumnya.
  \item Hitunglah $\frac{D_5}{5!}$ lalu bandingkan dengan
        $\frac1\eu \approx 0.3679$: peluang bahwa sebuah
        pesta besar yang dikocok tertukar sepenuhnya adalah
        $\frac1\eu$ --- penampilan singkat ketiga tetapan ini, sesudah
        undian dan sekretaris pada \cref{pb:g12:exp:1}.
        (Sebabnya: rumus pada pertanyaan 14 adalah awal sebuah
        deret masyhur untuk $\eu^{-1}$, yang diceritakan dalam jilid
        universitas.)
  \item Tukar kado di antara $10$ sahabat: namanya diambil
        secara acak seragam. Berapa peluang pengambilannya
        sah (tak seorang pun mengambil namanya sendiri), dan berapa kali
        pengambilan ulang yang patut mereka harapkan?
\end{enumerate}

\medskip
\noindent\textbf{Bagian IV --- Mencacah dua kali, menang dua kali.}
\begin{enumerate}[resume]
  \item Lema jabat tangan: pada pesta mana pun, menjumlahkan atas para tamunya
        banyaknya tangan yang dijabat masing-masing akan mencacah tiap jabat
        tangan tepat dua kali. Simpulkan bahwa \emph{banyaknya tamu
        yang menjabat sejumlah ganjil tangan selalu genap} ---
        lalu periksalah bahwa pernyataan itu masuk akal pada pesta bertiga.
  \item Buktikan permatanya,
        $1^3 + 2^3 + \dots + n^3 = (1 + 2 + \dots + n)^2$, dengan
        induksi, lalu periksalah untuk $n = 3$. (Jumlah Gauss kecil,
        yang dikuadratkan, ternyata mencacah pangkat tiga.)
  \item Identitas Vandermonde (\cref{exo:g12:comb:9}) lewat
        lintasan: tafsirkan $\binom{2n}{n}$ sebagai lintasan kekisi
        semacam pertanyaan 3 dari $(0,0)$ ke $(n,n)$, potonglah setiap
        lintasan di perpotongannya dengan diagonal lawannya, lalu jelaskan
        bagaimana $\sum_j \binom nj^2$ muncul.
  \item Penutup --- empat jurus sang pencacah, masing-masing satu baris dengan
        satu contoh dari soal ini: kalikan tahapannya dan jumlahkan
        kasusnya; sandikan dengan cerdik (bintang dan batang, kata lintasan);
        cacahlah hal yang sama dua kali (panitia-beserta-ketua,
        jabat tangan); kurangkan yang tak diinginkan lalu perbaiki
        kelebihan cacahnya (\omterm{def:g12:comb:permutation}{permutasi} kacau). Lalu perhatikan ke mana pencacahan
        ini bekerja selanjutnya: peluang, dan lintasan pada
        bab matriks dan graf.

\end{enumerate}
\end{problem}
